\begin{pastexam}[label={px:df3-114-ds-01}]{地方三等資訊處理 114 年 第一題}
（申論；20 分）

請回答下列問題：（每小題 4 分，共 20 分）
\begin{enumerate}[leftmargin=2.2em]
\item[㈠] 請說明堆疊（Stack）及佇列（Queue）那一種資料結構較適合用來進行後序（Postfix）運算式的計算？
\item[㈡] 請說明在二元搜尋樹中，前序（Preorder）走訪、中序（Inorder）走訪、後序（Postorder）走訪、層序
（Level-order）走訪那一種走訪順序可得到遞增的鍵值？
\item[㈢] 請說明在使用雜湊表時，若使用鏈結串列（chaining）處理碰撞（collision）問題，則搜尋的平均時間
複雜度為下列何者？O(1)、O(log n)、O(n)或 O(sqrt\{n\})。
\item[㈣] 請說明若一個圖 G（V, E）的頂點數|V|為 $n$，而邊數|E|接近 $n^2$，則相鄰串列（adjacency list）、
相鄰矩陣（adjacency matrix）、或邊列表（edge list）中，那一種資料結構最適合用來儲存該圖？
\item[㈤] 下列那幾項演算法可用於找出圖的最小生成樹（Minimum Spanning Tree）：Dijkstra 演算法、
Floyd-Warshall 演算法、Prim 演算法、Bellman-Ford 演算法？
\end{enumerate}
\end{pastexam}

\begin{pastsoln}
五個小題分屬資結書第~6、9、12、21、24~章，㈣ ㈤ 都是圖，這裡歸在第~21~章。

\textbf{㈠ 堆疊。} 後序運算式由左往右讀：讀到運算元就 push；讀到運算子就 pop 出兩個運算元（先 pop 的
是右運算元），算完把結果 push 回去；讀完時堆疊裡剩下的一個數就是答案（資結書第~6.3.2~節）。 運算子
要用的是「最近讀到、還沒用掉」的兩個值，正是後進先出。 例如 \texttt{2 3 + 4 *}：push 2、push 3，
讀到 + pop 出 3、2 算 5 push，push 4，讀到 * pop 出 4、5 算 20。 佇列是先進先出，取出的會是最早讀到的
運算元，順序不對。

\textbf{㈡ 中序走訪。} 二元搜尋樹每個節點的左子樹鍵值都比它小、右子樹都比它大。 中序走訪先走完整個
左子樹、再訪問節點、再走右子樹，由歸納法，每棵子樹的中序都是遞增的，接起來整棵也是遞增
（資結書定理~9.1）。 其他三種不行，例如根是 2、左子 1、右子 3 的樹：前序 2, 1, 3；後序 1, 3, 2；層序
2, 1, 3，都不是遞增。

\textbf{㈢ $O(1)$。} 鏈結法把雜湊到同一格的鍵串成一條串列。 表有 $m$ 格、存 $n$ 個鍵，負載因子
$\alpha = n/m$。 在雜湊函數把鍵均勻分到各格的假設下，搜尋的平均時間是 $\Theta(1 + \alpha)$（資結書
第~12.2.2~節、定理~12.2）：$1$ 是算雜湊值、找到那一格，$\alpha$ 是串列的平均長度。 實務上會讓表的大小
跟著 $n$ 調整（例如 $\alpha$ 超過某個常數就擴表重排），$\alpha = O(1)$，所以平均搜尋是 $O(1)$。 最差
情況（全部擠在同一格）才是 $O(n)$。

\textbf{㈣ 相鄰矩陣。} 邊數接近 $n^2$ 是稠密圖。 相鄰矩陣用 $\Theta(n^2)$ 空間，和邊數同一個數量級，
沒有浪費；而且判斷兩點是否相鄰、取得邊權都是 $O(1)$（資結書第~21.2~節）。 相鄰串列空間
$\Theta(n + |E|) = \Theta(n^2)$，數量級一樣，但每條邊多存一個指標，常數比較大，查兩點是否相鄰要掃
串列。 邊列表也是 $\Theta(|E|)$，但要找某個頂點的鄰居得掃過全部的邊，最不方便。

\textbf{㈤ 只有 Prim 演算法。} Prim 從一個頂點開始，每次加入連到樹外最輕的邊，得到最小生成樹
（資結書第~24.3~節）。 其他三個都是最短路徑演算法，解的是不同的問題：Dijkstra、Bellman-Ford 是單源
最短路徑，Floyd-Warshall 是所有點對最短路徑（資結書第~25~章）。 Dijkstra 的程式結構和 Prim 很像，
但 Dijkstra 比的是「從起點到這個頂點的距離」，Prim 比的是「這條邊本身的權重」，得到的最短路徑樹一般
不是最小生成樹。 例如三角形 $a$--$b$ 權 2、$b$--$c$ 權 2、$a$--$c$ 權 3，從 $a$ 出發的最短路徑樹是
\{$a$--$b$, $a$--$c$\}，總權重 5；最小生成樹是 \{$a$--$b$, $b$--$c$\}，總權重 4。
\end{pastsoln}
